\section{The complete graph}\label{sec:complete}\label{sec:completegraph}

On the complete graph every face ties at first order, and a few exact facts resolve the
tie. The results of this section and the next are those of our preprint~\cite{paperB}, and
Table~\ref{tab:overlap} says which of them follow from the general theory and which use the
exact arguments of the preprint. On a complete face the tree sum of Proposition~\ref{prop:constant}(b) has a closed form. It is
classical\ifarxivbuild, and we include a short proof\fi.

\begin{lemma}[weighted Cayley formula]\label{lem:weighted-cayley}
For $x_1,\ldots,x_k>0$ and the conductances $c_{ij}=x_ix_j$,
\[
\tree(c)=x_1x_2\cdots x_k\,(x_1+\cdots+x_k)^{k-2}.
\]
With every $x_i=1$ this is Cayley's count $k^{k-2}$ of the spanning trees of
the complete graph $K_k$.
\end{lemma}

\ifarxivbuild
\begin{proof}
Put $X=\sum_ix_i$. The Laplacian of the conductances is $\operatorname{diag}(Xx_i)-xx^\top$,
so without its first row and column it is
$\operatorname{diag}(Xx_2,\ldots,Xx_k)-x'x'^\top$ with $x'=(x_2,\ldots,x_k)$. The
matrix determinant lemma gives the determinant
$X^{k-1}\prod_{i\ge2}x_i\,(1-\sum_{i\ge2}x_i/X)=X^{k-2}\prod_ix_i$, and this is
$\tree(c)$ by the matrix-tree theorem~\cite{bollobas1998}. That theorem holds for real
conductances, since both sides are polynomials in the $c_{ij}$ that agree for integer
conductances, the multigraph case.
\end{proof}
\else
Its proof, by the matrix-tree theorem~\cite{bollobas1998} and the matrix determinant lemma, is in
Section~\ref{sec:complete} of the arXiv version of this paper~\cite{arxivC1}.
\fi

\begin{example}[the complete face]\label{ex:complete-face}
Let $F$ be a face of $k\ge2$ states with $q_{ij}=1$ for all $i\ne j$ in $F$, and
any exit rates, for example a face of $Q=J-dI$, where $q_i=d-1$. Then $F$ is
reversible with $\nu=\ones$, so by Proposition~\ref{prop:constant}(b) we may take
$r_\alpha=l_\alpha=\sqrt\alpha$ entrywise, and the conductances
$c_{ij}(\alpha)=\sqrt{\alpha_i\alpha_j}$ are products. Put
$A=\sum_{i\in F}\sqrt{\alpha_i}$. Lemma~\ref{lem:weighted-cayley} gives
$\tree(c(\alpha))=(\prod_i\alpha_i)^{1/2}A^{k-2}$, so for every
$\alpha\in\Delta_F^\circ$
\[
C_F(\alpha)=\frac{(\mu\cdot\sqrt\alpha)\,A^{k/2}}{(4\pi)^{(k-1)/2}\,(\prod_i\alpha_i)^{3/4}}
=\frac{(\mu\cdot\sqrt\alpha)\,A}{\sqrt{\det(2\pi\Sigma_\alpha)}},\qquad
J_F(\alpha)=\sum_{i\in F}q_i\alpha_i-(A^2-1).
\]
The second form of $C_F(\alpha)$ uses $\det\Sigma_\alpha=2^{k-1}(\prod_i\alpha_i)^2/\tree(c(\alpha))$.
For $Q=J-dI$ the rate is $J_F(\alpha)=d-A^2$, the rate of Dupuis and
Liu~\cite[Eq.~(3.2)]{dupuisliu2015}. At the center $\alpha=\ones/k$ we have $A^2=k$
and $C_F=\mu(F)k^{k-1/2}(4\pi)^{-(k-1)/2}$, which is Corollary~\ref{cor:symface}
with $\tree(q_F)=k^{k-2}$.
\end{example}

From now on in this section $Q=J-dI$. The chain switches to each other state with
probability $h$ and keeps its state with probability $p=1-(d-1)h$. The
following symbols of~\cite{paperB} are local to this section and to
Appendix~\ref{app:complete}. Put
\[
\sigma=p-h,\qquad u=\frac hp,\qquad v=\frac h\sigma=\frac u{1-u},\qquad
V_N=\frac{p^{N-1}}d .
\]
So $u$ is the odds of a switch to one given state against no switch, and $V_N$ is the
probability of a vertex $Ne_i$ for the uniform start, that is, $M_{\{i\}}$. For a composition
$n$ let $W_n(j)$ be the number of words with the counts $n$ and $j$ switches, and
$S_n(u)=\sum_jW_n(j)u^j$. For a
composition with some zero coordinates, $S_n$ is the polynomial of its positive coordinates,
which does not depend on $d$. By~\eqref{eq:runform}, the uniform start has
\begin{equation}\label{eq:SnVN}
P_N(n)=V_NS_n(u)\qquad(0<p\le1),
\end{equation}
and so does every start that is constant on $\supp n$, with $V_N$ replaced by
$\mu_ap^{N-1}$ for $a\in\supp n$. A composition is \emph{balanced} if its positive coordinates
differ by at most $1$. For the uniform start all faces of $k$ states have the same face maximum
$M_F$, and all balanced compositions with $k$ positive coordinates have the same probability. For
$p\ge1/d$ this probability is $M_F$, by Theorem~\ref{thm:simplex-balance} below. Put
\[
a_k=\frac12\log(4\pi)-\frac{k+\frac12}{k-1}\log k,\qquad
D_k=e^{-(k-1)a_k}=\frac{k^{k+1/2}}{(4\pi)^{(k-1)/2}}\qquad(k\ge2),
\]
and $D_1=1$. Numerically $a_2=-0.4674$, $a_3=-0.6571$ and $a_4=-0.8139$.

\begin{corollary}[the complete graph]\label{cor:paperB}
A face $F$ of $k$ states has $\lambda_F=d-k$ and $\alpha^*_F$ uniform, and if $\mu(F)>0$ then
$K_F=\frac{\mu(F)}kD_k$. So $\Phi_1(F)=d-1$ for every face, and all faces tie at $\tau=1$. With
the uniform start, the crossing of a face of $k\ge2$ states with a vertex has
$\Delta=\delta=k-1$ and $b_{FG}=a_k$.
\end{corollary}

\ifarxivbuild
\begin{proof}
Each $Q_F$ is symmetric with row sums $-(d-k)$, and the complete graph on $F$ has $k^{k-2}$
spanning trees by Lemma~\ref{lem:weighted-cayley}. So Corollary~\ref{cor:symface} gives
$\lambda_F$, $\alpha^*_F$ and
\[
K_F=\mu(F)k^{(k+1)/2}k^{(k-2)/2}(4\pi)^{-(k-1)/2}=\frac{\mu(F)}kD_k.
\]
A vertex has $K=\frac1d$ and $\lambda=d-1$ for the uniform start, so
$b_{FG}=-\frac1{k-1}\log(dK_F)=-\frac1{k-1}\log D_k=a_k$.
\end{proof}
\else
Its proof, from Corollary~\ref{cor:symface} and Lemma~\ref{lem:weighted-cayley}, is in
Section~\ref{sec:complete} of the arXiv version.
\fi

By Example~\ref{ex:complete-face} at the center, $D_k=dK_F=k/\sqrt{\det(2\pi\Sigma_k)}$ with
$\Sigma_k=\Sigma_{\ones/k}=(2/k^2)(I-J/k)$ in the coordinates $\alpha'$. So $D_k$ is the mass $k$
of a face, in units of a vertex, over the Gaussian volume of the fluctuations of its
fractions.

\begin{theorem}[balanced crossings]\label{thm:simplex-crossing}
Let $\mu$ be uniform and let $F$ be a face of $2\le k\le d$ states.
\begin{enumerate}[(i)]
\item For every $N\ge k$ there is exactly one $p=p_{N,k}\in(0,1)$ with $M_F=V_N$. The face
maximum $M_F$ is larger than $V_N$ below $p_{N,k}$ and smaller above it, and so is the
probability of a balanced composition of $F$. The crossing $p_{N,k}$ is strictly increasing in $N$,
$p_{k,k}=1/(1+(d-1)(k!)^{-1/(k-1)})$, and every crossing has $u\le(k!)^{-1/(k-1)}<1$, so
$p_{N,k}>1/d$.
\item For large $N$, uniformly for $1\le T\le N^{1/3}$,
\[
\log\frac{M_F}{V_N}=(k-1)\Bigl[T-L+\frac12\log T-a_k\Bigr]
+O\Bigl(\frac1T+\frac{T^2}N\Bigr).
\]
With $T=L-\frac12\log L+t$, the ratio $M_F/V_N$ tends to $e^{(k-1)(t-a_k)}$, uniformly for
$t$ in compact sets.
\item The value $T_{N,k}=N(1-p_{N,k})/(d-1)$ of $T$ at the crossing satisfies
\[
T_{N,k}+\frac12\log T_{N,k}=L+a_k+O(1/L),
\]
and
\begin{equation}\label{eq:simplex-root-expansion}
T_{N,k}=L-\frac12\log L+a_k+\frac{\log L}{4L}+O\Bigl(\frac1L\Bigr).
\end{equation}
\item $a_2>a_3>a_4>\cdots$.
\end{enumerate}
\end{theorem}

For $d=2$ the chain gives the directions of the persistent walk, and $T_{N,2}$ is the
crossing of its middle compositions and its endpoints, with $a_2=\frac12\log(\pi/8)$. The proof
is in Appendix~\ref{app:complete}\arxivpointer.

\begin{theorem}[exact balancing]\label{thm:simplex-balance}
Let $1/d\le p<1$, and let $\mu$ be constant and positive on a face $F$. Moving one unit from a
coordinate $b$ of a composition with support $F$ to another coordinate $a$, where $n_b\ge n_a+2$ and
$n_a\ge1$, strictly increases the probability of the composition. So the most likely compositions with support
$F$ are exactly the balanced ones, and they all have the same probability.
\end{theorem}

For an integer $n\ge1$ put
\begin{equation}\label{eq:refresh-polynomial}
B_n(w)=\sum_{m=1}^n\binom{n-1}{m-1}\frac{w^m}{m!}
=\frac wn L^{(1)}_{n-1}(-w),
\end{equation}
where $L^{(1)}_{n-1}$ is the Laguerre polynomial~\cite[Eq.~18.5.12]{dlmf}. The proof
in Appendix~\ref{app:complete}\arxivpointer{} sums over the numbers of runs first. This writes the law as an
integral of a product of one polynomial $B_{n_i}$ per state, and $B_n(w)$ is log-concave in $n$
for $w>0$. The condition
$p\ge1/d$ holds for $T\le N/d$, and so at every crossing by
Theorem~\ref{thm:simplex-crossing}(i).

\begin{corollary}[exactly one mode change]\label{cor:simplex-modes}
Let $\mu$ be uniform, and fix $d\ge2$ and $C>1$. There is $N_0$ such that for $N\ge N_0$ and
every $p$ with $0<T\le CL$ the following hold.
\begin{enumerate}[(i)]
\item If $p>p_{N,d}$, the global modes are exactly the $d$ vertices.
\item If $p<p_{N,d}$, they are exactly the balanced compositions with $d$ positive
coordinates.
\item If $p=p_{N,d}$, they are these two classes and nothing else.
\end{enumerate}
Moreover $p_{N,2}<p_{N,3}<\cdots<p_{N,d}$ and
$p_{N,d}-p_{N,k}=(d-1)(a_k-a_d+o(1))/N$. So all the crossings of balanced compositions with the
vertices lie in a $p$-interval of length $(d-1)(a_2-a_d+o(1))/N$, and as $p$ decreases the
balanced compositions with $d$ positive coordinates cross first.
\end{corollary}

On the triangle the middle composition
$(\frac N3,\frac N3,\frac N3)$ overtakes the vertices while the middles of the edges are still
below them by the factor $e^{-(a_2-a_3)}=0.827$. The proof is in Appendix~\ref{app:complete}\arxivpointer.
In the tie window at $\tau^*=1$ this is Theorem~\ref{thm:window}, whose envelope $E$ has a
single kink, at $t=a_d$.

\begin{remark}[an intermediate face can be a mode at finite $N$]\label{rem:simplex-edge-mode}
Corollary~\ref{cor:simplex-modes} does not hold for every $N$. Take $d=3$, $N=4$ and
$p=10/17$, so $u=7/20$ and $0<T=14/17<\log4$. Counting words gives
$P_4(2,2,0)/V_4=2u+2u^2+2u^3=4123/4000>1$ and $P_4(2,1,1)/V_4=6u^2+6u^3=3969/4000<1$.
Since $p>1/3$, Theorem~\ref{thm:simplex-balance} shows that every global mode is a vertex or
a balanced composition, so the three edge centers are the only global modes. Hence $N_0>4$
when $d=3$, for every $C>1$.
\end{remark}

